\clearpage\SourceLeaf{165}%
\pdfbookmark[1]{Summatio quarumdam serierum singularium. Besprechung 1808 September 19}{notice-1808-09-19}
\begin{center}
\rule{.98\linewidth}{.3pt}\par
Göttingische Gelehrte Anzeigen.\quad 1808 September 19.\par
\rule{.98\linewidth}{.3pt}
\end{center}

Eine vom Herrn Professor \textsc{Gauss} der Königlichen Gesellschaft der Wissenschaften vorgelegte Vorlesung,
\begin{center}\emph{Summatio quarumdam serierum singularium},\end{center}
soll eine bemerkenswerte Untersuchung zur Kreisteilung, deren Grundlage bereits in den \emph{Disquisitionibus Arithmeticis} gelegt war, ausführlicher und in größerer Allgemeinheit entwickeln, sie mit vollständigen Beweisen versehen und ihren unerwarteten Zusammenhang mit anderen wichtigen Wahrheiten zeigen. Es bezeichne $n$ eine Primzahl, $k$ eine beliebige ganze Zahl, die durch $n$ nicht teilbar ist, und $\omega$ den Bogen $\frac1n360^\circ$. Die verschiedenen, unter den Zahlen $1,2,3,4,\dotsfour n-1$ vorkommenden quadratischen Reste von $n$ seien mit $a,a',a''$ usw. bezeichnet; die nach deren Ausschluss verbleibenden Zahlen, also die quadratischen Nichtreste von $n$, dagegen mit $b,b',b''$ usw. Dann ist im angeführten Werk in Art.~356 bewiesen, dass in dem Fall, in dem $n$ die Form $4n+1$ hat,
\[
\left.\begin{aligned}
&\cos ak\omega+\cos a'k\omega+\cos a''k\omega+\text{usw.}\\
&-\cos bk\omega-\cos b'k\omega-\cos b''k\omega-\text{usw.}
\end{aligned}\right\}=\pm\surd n.
\]
und
\[
\left.\begin{aligned}
&\sin ak\omega+\sin a'k\omega+\sin a''k\omega+\text{usw.}\\
&-\sin bk\omega-\sin b'k\omega-\sin b''k\omega-\text{usw.}
\end{aligned}\right\}=0
\]
\SourceLeaf{166}%
gilt, während in dem Fall, in dem $n$ die Form $4m+3$ hat, die Summe der ersten Reihe $=0$ und die der zweiten $=\pm\surd n$ wird. Das Vorzeichen der Wurzelgröße hängt vom Wert der Zahl $k$ oder vielmehr von ihrer Beziehung zu $n$ ab und lässt sich für \emph{alle} Werte von $k$ bei einem gegebenen Wert von $n$ leicht bestimmen, sobald es für \emph{einen} bestimmt ist. Man kann nämlich zeigen, dass für alle Werte von $k$, die quadratische Reste von $n$ sind, stets \emph{dasselbe} Vorzeichen gilt und dann das entgegengesetzte für alle quadratischen Nichtreste von $n$. Da die Untersuchung im angeführten Werk bereits so weit geführt war und nur noch die Bestimmung des Vorzeichens für irgendeinen Wert von $k$ fehlte, hätte man glauben sollen, dass sich nach Erledigung der Hauptsache diese genauere Bestimmung leicht ergänzen ließe. Das lag umso näher, als die Untersuchung von Zahlenbeispielen sofort ein äußerst einfaches Ergebnis liefert: Für $k=1$ oder für alle Werte, die quadratische Reste von $n$ sind, muss die Wurzelgröße in den obigen Formeln stets \emph{positiv} genommen werden. Doch bei der Suche nach einem Beweis dieser Beobachtung stoßen wir auf ganz unerwartete Schwierigkeiten. Das Verfahren, das so befriedigend zur Bestimmung des absoluten Werts jener Reihen führte, erweist sich als völlig unzureichend, wenn die Vorzeichen vollständig bestimmt werden sollen. Den \emph{metaphysischen} Grund dieses Phänomens (um den bei französischen Mathematikern üblichen Ausdruck zu gebrauchen) muss man darin suchen, dass die Analysis bei der Kreisteilung zwischen den Bögen $\omega,2\omega,3\omega\ldots(n-1)\omega$ keinen Unterschied macht, sondern sie alle auf dieselbe Weise erfasst. Da die Untersuchung dadurch ein neues Interesse erhält, fand Herr Professor \textsc{Gauss} darin gewissermaßen eine Aufforderung, nichts unversucht zu lassen, um die Schwierigkeit zu überwinden. Erst nach vielen und verschiedenartigen vergeblichen Versuchen gelang ihm dies auf einem auch für sich bemerkenswerten Weg. Er geht nämlich von der Summation einiger Reihen aus, deren Glieder die folgende Form haben:
\[
\frac{(1-x^m)(1-x^{m-1})(1-x^{m-2})\ldots(1-x^{m-\mu+1})}{(1-x)\ (1-xx)\ (1-x^3)\dotsfour(1-x^\mu)}
\]
Bezeichnen wir der Kürze halber eine solche Funktion mit $(m,\mu)$, die, wie in der Abhandlung gezeigt wird, stets ein \emph{Polynom} in $x$ ist, so brechen die Reihen
\SourceLeaf{167}%
\[
1-(m,1)+(m,2)-(m,3)+\text{usw.}
\]
\[
1+x^{\frac12}(m,1)+x(m,2)+x^{\frac32}(m,3)+\text{usw.}
\]
nach dem Glied an Stelle $m+1$ ab, sofern $m$ eine positive ganze Zahl bezeichnet. Für gerade Werte von $m$ wird die Summe der ersten Reihe
\[
=(1-x)(1-x^3)(1-x^5)\ldots(1-x^{m-1})
\]
und $=0$ für ungerade Werte von $m$; die Summe der zweiten Reihe wird dagegen stets
\[
=(1+x^{\frac12})(1+x)(1+x^{\frac32})\ldots(1+x^{\frac12m})
\]
Auch für gebrochene und negative Werte von $m$ führt die Summation dieser Reihen zu interessanten Ergebnissen, obwohl diese für den gegenwärtigen Zweck nicht nötig sind. Wir begnügen uns damit, hier nur eines davon anzugeben. Die unendliche Reihe
\[
1+x+x^3+x^6+x^{10}+\text{usw.}
\]
deren Exponenten die Dreieckszahlen sind, ist das Produkt aus den Faktoren
\[
\frac{1-xx}{1-x}\times\frac{1-x^4}{1-x^3}\times\frac{1-x^6}{1-x^5}\times\frac{1-x^8}{1-x^7}\ \text{usw.}
\]
\begin{samepage}
oder, wenn man es vorzieht, aus
\[
(1+x)^2(1+xx)^2(1+x^3)^2(1+x^4)^2\ \text{usw.}
\]
multipliziert mit
\[
(1-x)(1-xx)1-x^3)(1-x^4)\ \text{usw.}
\]
\end{samepage}

Darzulegen, wie diese Summationen auf den Hauptgegenstand angewandt werden, würde uns hier zu weit führen. Wir dürfen die Leser umso eher auf die Abhandlung selbst verweisen, als sie bald im Druck erscheinen wird. Die oben angeführten Summationen sind nur ein spezieller Fall der Summation der folgenden Reihen:
\[
1+\cos k\omega+\cos4k\omega+\cos9k\omega+\text{usw.}+\cos(n-1)^2k\omega=T
\]
\[
\sin k\omega+\sin4k\omega+\sin9k\omega+\text{usw.}+\sin(n-1)^2k\omega=U
\]
die in der Abhandlung für alle Werte von $k$ und ohne die Einschränkung
\SourceLeaf{168}%
behandelt wird, dass $n$ eine Primzahl sein muss. Es wird nämlich gezeigt, dass
\[
T=\pm\surd n,\quad T=\pm\surd n,\quad T=0,\quad T=0
\]
und
\[
U=\pm\surd n,\quad U=0,\quad U=0,\quad U=\pm\surd n
\]
gilt, je nachdem $n$ die Formen $4m,4m+1,4m+2,4m+3$ der Reihe nach hat. Das Vorzeichen der Wurzelgröße hängt hier wiederum von $k$ ab. Seine Bestimmung, die die Unterscheidung vieler einzelner Fälle erfordert, wird auf zwei verschiedenen Wegen so entwickelt und bewiesen, dass nichts zu wünschen übrig bleiben wird. Der Vergleich dieser beiden Wege führt außerdem zu dem folgenden sehr bemerkenswerten Satz: Ist $n$ das Produkt aus beliebig vielen verschiedenen ungeraden Primzahlen $a,b,c,d$ usw., von denen insgesamt $\mu$ die Form $4m+3$ haben, und kommen unter jenen Faktoren insgesamt $\nu$ vor, von denen jeder das Produkt der übrigen (also der Reihe nach $\frac na,\frac nb,\frac nc,\frac nd$ usw.) zum quadratischen Nichtrest hat, so ist $\nu$ gerade, wenn $\mu$ die Form $4m$ oder $4m+1$ hat, dagegen ungerade, wenn $\mu$ die Form $4m+2$ oder $4m+3$ hat. Der bekannte Fundamentalsatz über die quadratischen Reste ist nur ein spezieller Fall dieses Satzes; umgekehrt lässt sich der eben angegebene Satz leicht aus diesem Fundamentalsatz ableiten. Durch diese Untersuchungen besitzt man also zugleich einen vierten Beweis dieses wichtigen Satzes, den der Verfasser zuerst auf zwei ganz verschiedenen Wegen in den \emph{Disquisitionibus Arithmeticis} und vor Kurzem auf einem dritten, ebenso verschiedenen Weg in einer eigenen Abhandlung bewiesen hatte.
\BookRule{.28\linewidth}
